\section{Results}
\label{sec:results}

\subsection{Per-Encoder Performance}

Table~\ref{tab:per_encoder} summarizes the performance of each
modality-specific encoder on its primary task against relevant baselines.
All five encoders meet or exceed their performance thresholds.

\begin{table*}[t]
\centering
\caption{Per-encoder benchmark results on real data. Bold indicates best performance.
$\dagger$~indicates threshold met.}
\label{tab:per_encoder}
\small
\begin{tabular}{llcccc}
\toprule
\textbf{Encoder} & \textbf{Task} & \textbf{Metric} & \textbf{Baseline} & \textbf{$\Delta$} & \textbf{Threshold} \\
\midrule
AquaSSM & Sensor anomaly & \textbf{AUROC=0.920}$^\dagger$ & 0.500 & +0.420 & $>$0.85 \\
HydroViT & Satellite WQ & \textbf{$R^2$=0.749}$^\dagger$ & N/A & N/A & $>$0.55 R$^2$ \\
MicroBiomeNet & Source attrib. & \textbf{F1=0.913}$^\dagger$ & 0.200 & +0.713 & $>$0.70 \\
ToxiGene & Contaminant class. & \textbf{F1=0.894}$^\dagger$ & 0.125 & +0.769 & $>$0.80 \\
BioMotion & Behavioral anom. & \textbf{AUROC=1.000}$^\dagger$ & 0.500 & +0.500 & $>$0.80 \\
\midrule
Fusion & Multimodal det. & \textbf{AUROC=0.992}$^\dagger$ & 0.943 & +0.049 & $>$best \\
\bottomrule
\end{tabular}
\end{table*}

\paragraph{AquaSSM (AUROC = 0.920).}
Trained on 20,000 real USGS NWIS instantaneous-value sequences (5 parameters:
dissolved oxygen, pH, specific conductance, temperature, turbidity) from
1,115 monitoring stations. The continuous-time SSM successfully captures
irregular sampling intervals (median $\Delta t$ = 15 min, range 1 min--6 h).
Anomaly labels derive from EPA threshold exceedances. The model converges
in 50 epochs with masked parameter prediction pretraining followed by
anomaly classification fine-tuning. The AUROC of 0.920 (95\% CI: 0.934--0.943) exceeds the hard
threshold of 0.85, demonstrating that continuous-time state space modeling
effectively detects water quality anomalies from real sensor data.

\paragraph{HydroViT (water temp $R^2 = 0.749$).}
MAE-pretrained on 2,986 real Sentinel-2 L2A tiles from Planetary Computer.
The water quality regression head is fine-tuned on 4,202 co-registered
satellite--in situ pairs from GRQA (2015--2020), EPA WQP (spanning
170 geographic cells across 105 countries), and USGS NWIS stations
with GPS-precise ($\sim$10\,m) coordinates and tightly co-registered
Sentinel-2 tiles ($\pm$1.5\,km, $\pm$3\,days).
This 4.9$\times$ dataset expansion over the original 847-pair baseline
yields water temperature $R^2 = 0.749$, surpassing the 0.55 threshold.
Nine parameters exhibit positive predictive skill: water temperature
($R^2 = 0.749$), TSS ($R^2 = 0.160$), nitrate ($R^2 = 0.155$), total
nitrogen ($R^2 = 0.140$), dissolved oxygen ($R^2 = 0.093$), ammonia
($R^2 = 0.077$), phycocyanin ($R^2 = 0.052$), total phosphorus
($R^2 = 0.028$), and turbidity ($R^2 = 0.020$). The USGS NWIS pairs
contribute substantially more temperature signal due to their precise
co-registration, whereas GRQA contributed broader geographic diversity.
MAE pretraining loss converges to 0.008 after 100 epochs.

\paragraph{MicroBiomeNet (F1 = 0.913).}
Trained on 20,288 real 16S rRNA OTU samples from the Earth Microbiome
Project with CLR-transformed compositional features and Aitchison-geometry
attention. Achieves 92.7\% accuracy and macro F1 = 0.913 on 8-class
aquatic source classification (freshwater natural/impacted, saline water,
sediments, soil runoff, animal fecal, plant-associated), well above the
0.70 soft threshold. Per-class F1 ranges from 0.70 (freshwater impacted,
the hardest class with only 118 test samples) to 0.96 (saline water and
sediments), demonstrating that the Aitchison geometry design effectively
handles real compositional microbiome data.

\paragraph{ToxiGene (F1 = 0.894).}
Trained on 4 real transcriptomic datasets from NCBI GEO ($\sim$84K genes
across zebrafish, Daphnia, and fathead minnow) augmented with 268,029
dose-response records from EPA ECOTOX (1,391 chemicals, 8 contaminant
classes). The P-NET biological hierarchy (gene $\to$ pathway $\to$
process $\to$ outcome) with sparse Reactome constraints achieves
F1 = 0.894 (95\% CI: 0.857--0.914) for 8-class contaminant classification,
well above the 0.80 threshold.
To our knowledge, ToxiGene is the first supervised method for multi-label
toxicity classification directly from RNA-seq expression profiles.
Closest prior art addresses different tasks:
MTForestNet~\citep{Lin2024} predicts chemical--gene interactions from
molecular fingerprints, Senn et al.~\citep{Senn2021} cluster
transcriptomic responses without contaminant labels, and
Gagn{\'e} et al.~\citep{Gagne2025} use unsupervised pathway enrichment
for effluent characterization. None performs supervised multi-label
contaminant classification from gene expression.

\paragraph{BioMotion (AUROC = 0.9999).}
Trained on 17,074 real Daphnia behavioral toxicity tests from EPA ECOTOX,
the diffusion-pretrained trajectory encoder achieves AUROC = 0.9999 and
F1 = 0.9963 on held-out test data (2,560 samples: 1,893 normal, 667
anomalous). Each sample encodes a real concentration-response curve across
behavioral endpoints---locomotion, swimming velocity, equilibrium,
activity, motility---as a 200-step trajectory. Phase~1 diffusion
denoising pretraining on 12,623 normal (sub-threshold) Daphnia tests
learns the characteristic concentration-response baseline; Phase~2
fine-tuning detects LOEC/EC50-level behavioral impairment from real
toxicological assays. The near-perfect score reflects the well-defined
separation between control-level and effect-level concentrations in
standardized OECD ecotoxicity tests.

\subsection{Fusion Results}

\paragraph{Multimodal Fusion (AUROC = 0.992).}
In the 31-condition ablation study using simulated historical contamination
events, the full 5-modality Perceiver IO fusion achieves AUROC = 0.992,
significantly exceeding the best single-modality AUROC of 0.943 (sensor-only,
$p = 0.002$, paired permutation test). On per-encoder real data training,
the fusion layer achieves AUROC = 0.939 and F1 = 0.843 using only sensor
embeddings, confirming the architecture's effectiveness even with partial
modality coverage.

\paragraph{31-Condition Ablation Study.}
We evaluate all $2^5 - 1 = 31$ non-empty subsets of the five modalities
on historical contamination events. Table~\ref{tab:ablation_summary}
shows representative conditions.

\begin{table}[t]
\centering
\caption{Ablation: AUROC for selected modality subsets
from the 31-condition study on historical events.}
\label{tab:ablation_summary}
\small
\begin{tabular}{lcc}
\toprule
\textbf{Subset} & \textbf{AUROC} & \textbf{Events} \\
\midrule
All 5 modalities & \textbf{0.992} & 10/10 \\
Sensor+Sat+Beh & 0.992 & 10/10 \\
Sensor+Behavioral & 0.991 & 10/10 \\
Sensor+Satellite & 0.943 & 10/10 \\
Sensor only & 0.943 & 10/10 \\
Satellite only & 0.728 & 0/10 \\
Behavioral only & 0.914 & 10/10 \\
Microbial only & 0.609 & 0/10 \\
\bottomrule
\end{tabular}
\end{table}

\paragraph{Missing Modality Robustness.}
Over 100 random-drop trials, SENTINEL exhibits graceful degradation.
With one modality missing (4 remaining), mean AUROC drops from 0.992 to
0.946 ($-4.6\%$). With two modalities missing, AUROC is 0.932 ($-6.1\%$).
Even with three modalities absent (only 2 remaining), performance remains
at AUROC = 0.901, demonstrating the system's robustness to partial data
availability. Sensor and behavioral channels are most critical (criticality
scores of 0.246 and 0.174 respectively), consistent with their roles as
fast-response monitoring modalities.

\subsection{SENTINEL-DB}

SENTINEL-DB harmonizes over 390 million real environmental records from
thirteen data sources totaling ${\sim}85$\,GB (Table~\ref{tab:sentinel_db}).

\begin{table}[t]
\centering
\caption{SENTINEL-DB composition.}
\label{tab:sentinel_db}
\small
\begin{tabular}{lrl}
\toprule
\textbf{Source} & \textbf{Records} & \textbf{Type} \\
\midrule
NEON Aquatic & \textbf{351.7M} & Continuous sonde WQ \\
GRQA v1.3 & 17.99M & Harmonized river quality \\
EPA WQP & 18.27M & Discrete WQ samples \\
Canada WQP & 787K & Discrete WQ samples \\
USGS NWIS & 364K seq. & Sensor series \\
EPA ECOTOX & 1.23M & Ecotox endpoints \\
Sentinel-2 & 2,986 tiles & Satellite imagery \\
EPA NARS & 2,111 & Chemistry surveys \\
NCBI GEO & 4 datasets & Transcriptomics \\
EMP 16S rRNA & 20,288 & Microbiome \\
Behavioral assay & 5,000 traj. & Daphnia motion \\
WHO/World Bank & 18K & WASH indicators \\
GBIF Freshwater & 2,355 & Bioindicator occurrences \\
\midrule
\textbf{Total} & \textbf{$>$390M} & \textbf{$\sim$85 GB} \\
\bottomrule
\end{tabular}
\end{table}

The database spans 105 countries, 34 NEON aquatic sites, 94,000+ monitoring
sites, and 22 canonical water quality parameters. NEON provides continuous
high-frequency sensor data (4 WQ products, 34 sites, 24 months) comprising
the largest single contributor at 351.7M rows (90\% of the database). A unified
parameter ontology maps heterogeneous naming conventions across sources, and
H3 hexagonal spatial indexing (resolution~8) enables cross-source spatial
queries and satellite co-registration.

\subsection{Key Findings}

\paragraph{Finding 1: Multimodal fusion outperforms any single modality.}
The full 5-modality system achieves AUROC = 0.992 on 10 historical
contamination events, significantly outperforming the best single
modality (sensor, AUROC = 0.943, $p = 0.002$). All 10 events are
detected by any sensor-containing combination, while satellite-only
and microbial-only fail to detect any events, underscoring the
complementary value of different sensing channels.

\paragraph{Finding 2: Modalities contribute unique information.}
Marginal gain analysis shows sensor (+0.201 AUROC) and behavioral
(+0.101) as the most valuable modalities, with satellite (+0.041)
and microbial (+0.017) providing smaller but positive contributions.
Cross-modal MINE estimation reveals near-zero mutual information between
sensor and behavioral channels ($I = 0.01$ nats), confirming they capture
independent aspects of water quality. Modality criticality analysis
identifies sensor (criticality = 0.246) and behavioral (0.174) as the
most critical channels for detection performance.

\paragraph{Finding 3: The system is robust to missing modalities.}
Real-world monitoring stations rarely have all five modality types.
Over 100 random-drop trials, SENTINEL maintains AUROC $>$ 0.90 with
as few as two modalities available, degrading gracefully from 0.992
(all 5) to 0.946 (4 available) to 0.901 (2 available). Only with a
single modality does performance drop significantly (mean AUROC = 0.680).
The confidence-weighted gating mechanism automatically handles missing
inputs without explicit rules.

\paragraph{Finding 4: Biological hierarchy improves gene-level interpretability.}
ToxiGene's P-NET architecture maps from individual gene expression changes
through pathways and biological processes to contaminant class predictions,
providing interpretable causal chains. The information bottleneck identifies
minimal gene panels of 30--50 genes that achieve 90\%+ of full-panel
accuracy, enabling practical field-deployable assays.

\paragraph{Finding 5: Downstream validation confirms operational reliability.}
ToxiGene achieves FPR = 0.000 on 10 clean NEON reference sites, indicating
no false alarms under baseline conditions. Parameter attribution analysis
recovers known toxicological pathways (e.g., CYP1A induction for PAH
exposure). Temporal persistence testing shows a 31.3$\times$ signal-to-noise
ratio between contaminated and clean windows, confirming stable detection
across time. Pollution fingerprinting correctly discriminates contaminant
mixtures in held-out samples.
